\documentclass[11pt,a4paper,openany,fontset=fandol]{ctexbook}
\usepackage[top=24mm,bottom=23mm,left=24mm,right=24mm,headheight=15pt,headsep=9mm]{geometry}
\usepackage{amsmath,amssymb,mathtools,bm}
\usepackage{booktabs,tabularx,longtable,array,multirow}
\usepackage{enumitem}
\usepackage{xcolor}
\usepackage[most]{tcolorbox}
\usepackage{tikz}
\usetikzlibrary{arrows.meta,positioning,calc,fit,shapes.geometric,decorations.pathreplacing}
\usepackage{pgfplots}
\pgfplotsset{compat=1.18}
\usepackage{fancyhdr}
\usepackage{titlesec}
\usepackage{hyperref}
\usepackage{microtype}
\usepackage{caption}
\usepackage{float}

\definecolor{titleblue}{HTML}{1B527C}
\definecolor{deepgreen}{HTML}{2E6D57}
\definecolor{softgreen}{HTML}{F1F8F4}
\definecolor{deepred}{HTML}{984543}
\definecolor{softred}{HTML}{FCF2F1}
\definecolor{gold}{HTML}{B4862A}
\definecolor{softgold}{HTML}{FFF9EA}
\definecolor{deepblue}{HTML}{315D83}
\definecolor{softblue}{HTML}{EFF6FC}
\definecolor{softgray}{HTML}{F5F6F7}
\definecolor{violet}{HTML}{68548B}
\definecolor{softviolet}{HTML}{F4F0FA}

\hypersetup{colorlinks=true,linkcolor=titleblue,urlcolor=titleblue,citecolor=deepgreen,pdfauthor={OpenAI},pdftitle={二维材料光电器件、人工突触与光谱实验：第十六至第二十章}}
\setlength{\parindent}{2em}
\setlength{\parskip}{0.25em}
\linespread{1.16}
\setlist[itemize]{leftmargin=2.2em,itemsep=0.35em,topsep=0.35em}
\setlist[enumerate]{leftmargin=2.5em,itemsep=0.35em,topsep=0.35em}
\renewcommand{\arraystretch}{1.38}

\titleformat{\chapter}[display]
  {\bfseries\color{titleblue}}
  {\Large 第\chinese{chapter}章}
  {0.45em}
  {\Huge}
\titlespacing*{\chapter}{0pt}{-12pt}{1.4em}
\titleformat{\section}{\Large\bfseries\color{titleblue}}{\thesection}{0.8em}{}
\titleformat{\subsection}{\large\bfseries\color{titleblue}}{\thesubsection}{0.7em}{}
\titlespacing*{\section}{0pt}{2em}{0.8em}

\pagestyle{fancy}
\fancyhf{}
\fancyhead[L]{\small\color{gray}二维材料光电器件、人工突触与光谱实验}
\fancyhead[R]{\small\color{gray}\leftmark}
\fancyfoot[C]{\thepage}
\renewcommand{\headrulewidth}{0.4pt}
\renewcommand{\headrule}{\hbox to\headwidth{\color{gray!28}\leaders\hrule height \headrulewidth\hfill}}
\renewcommand{\chaptermark}[1]{\markboth{第\chinese{chapter}章\quad #1}{}}

\newtcolorbox{mainline}[1][]{enhanced,breakable,colback=softgreen,colframe=deepgreen,colbacktitle=deepgreen,coltitle=white,fonttitle=\bfseries,title={本节主线},boxrule=0.7pt,arc=2mm,left=4mm,right=4mm,top=2mm,bottom=2mm,#1}
\newtcolorbox{translationbox}[1][]{enhanced,breakable,colback=softgreen,colframe=deepgreen,colbacktitle=deepgreen,coltitle=white,fonttitle=\bfseries,title={把术语翻译成实验对象},boxrule=0.7pt,arc=2mm,left=4mm,right=4mm,top=2mm,bottom=2mm,#1}
\newtcolorbox{warningbox}[1][]{enhanced,breakable,colback=softred,colframe=deepred,colbacktitle=deepred,coltitle=white,fonttitle=\bfseries,title={常见误判},boxrule=0.7pt,arc=2mm,left=4mm,right=4mm,top=2mm,bottom=2mm,#1}
\newtcolorbox{examplebox}[1][]{enhanced,breakable,colback=softgold,colframe=gold,boxrule=0.8pt,arc=2mm,left=4mm,right=4mm,top=3mm,bottom=3mm,#1}
\newtcolorbox{definitionbox}[1][]{enhanced,breakable,colback=softblue,colframe=deepblue,boxrule=0.7pt,arc=2mm,left=4mm,right=4mm,top=3mm,bottom=3mm,#1}
\newtcolorbox{auditbox}[1][]{enhanced,breakable,colback=softviolet,colframe=violet,colbacktitle=violet,coltitle=white,fonttitle=\bfseries,title={开放论文审计},boxrule=0.7pt,arc=2mm,left=4mm,right=4mm,top=2mm,bottom=2mm,#1}
\newtcolorbox{checkbox}[1][]{enhanced,breakable,colback=softgray,colframe=gray!45,colbacktitle=titleblue,coltitle=white,fonttitle=\bfseries,title={本章课题自检},boxrule=0.6pt,arc=2mm,left=4mm,right=4mm,top=2mm,bottom=2mm,#1}
\newtcolorbox{conclusionbox}[1][]{enhanced,breakable,colback=softblue,colframe=titleblue,colbacktitle=titleblue,coltitle=white,fonttitle=\bfseries,title={本章结论},boxrule=0.7pt,arc=2mm,left=4mm,right=4mm,top=2mm,bottom=2mm,#1}

\newcommand{\tickitem}{\item[$\square$]}
\newcommand{\dd}{\mathrm{d}}
\newcommand{\e}{\mathrm{e}}
\newcommand{\figcaption}[1]{\par\smallskip\noindent\textbf{图 \thefigure}\quad #1\par\medskip\stepcounter{figure}}
\setcounter{chapter}{15}
\setcounter{page}{118}
\setcounter{figure}{0}

\begin{document}

\chapter{从器件曲线到网络性能}

一条平滑的增强--抑制曲线只是网络映射的入口，不是识别率、能耗或吞吐量的直接证明。网络真正接收的是一组带有限动态范围、非线性、随机性和空间差异的可编程电导；外围电路还会继续引入量化、线阻、读噪声和校准成本。本章的任务，是把“器件曲线漂亮”逐步翻译为“网络中哪些误差可计算、可训练、可补偿”。

\begin{mainline}
器件到网络至少经过四次映射：
\[
\text{脉冲序列}\rightarrow G(n)\rightarrow w(G)\rightarrow \mathbf y=f(\mathbf W,\mathbf x)\rightarrow \text{任务指标}.
\]
每一步都必须给出定义、误差来源和实验边界。缺少其中任何一步，软件准确率都可能只是对平均曲线的演示。
\end{mainline}

\section{网络需要的不是“开关比”，而是可用权重窗口}

设器件可编程电导位于 $[G_{\min},G_{\max}]$。最常用的归一化方式是
\[
 g=\frac{G-G_{\min}}{G_{\max}-G_{\min}},\qquad 0\le g\le 1.
\]
但这一定义默认所有状态都能可靠写入和区分。若相邻状态间距小于读噪声和漂移的组合标准差，名义状态数不会转化为有效状态数。可以用保守判据
\[
\Delta G_k > q\sqrt{\sigma_{\mathrm{read},k}^2+\sigma_{\mathrm{drift},k}^2+\sigma_{\mathrm{cycle},k}^2}
\]
判断第 $k$ 个相邻状态是否可分；$q=3$ 对应约三倍标准差的间隔。

\begin{definitionbox}
\textbf{有效权重状态数。} 在给定读偏压、保持时间和容许误码率下，能够被独立写入并可靠区分的电导状态数量。它必须和测量带宽、采样时间及判据一起报告，不能由一条连续曲线目测得到。
\end{definitionbox}

\section{增强--抑制非线性：先拟合，再进入网络}

常见经验模型可写为
\begin{align*}
G_{\mathrm P}(n)&=G_{\min}+\Delta G\frac{1-\exp(-n/A_{\mathrm P})}{1-\exp(-N/A_{\mathrm P})},\\
G_{\mathrm D}(n)&=G_{\max}-\Delta G\frac{1-\exp[-(N-n)/A_{\mathrm D}]}{1-\exp(-N/A_{\mathrm D})}.
\end{align*}
$A_{\mathrm P}$ 与 $A_{\mathrm D}$ 控制曲率，而不是直接等于某个微观时间常数。若只给平均曲线而不报告重复循环，网络仿真会低估写入随机性。

\begin{figure}[H]
\centering
\begin{tikzpicture}
\begin{axis}[width=0.84\textwidth,height=6.2cm,xlabel={归一化脉冲序号 $n/N$},ylabel={归一化电导 $g$},xmin=0,xmax=1,ymin=0,ymax=1.05,grid=both,legend style={at={(0.03,0.97)},anchor=north west,draw=none,fill=white},samples=120]
\addplot[very thick,titleblue,domain=0:1]{(1-exp(-4*x))/(1-exp(-4))};
\addlegendentry{强非线性增强}
\addplot[very thick,deepgreen,dashed,domain=0:1]{x};
\addlegendentry{理想线性}
\addplot[very thick,gold,domain=0:1]{1-(1-exp(-4*(1-x)))/(1-exp(-4))};
\addlegendentry{不对称抑制}
\end{axis}
\end{tikzpicture}
\caption{教学模拟：相同动态范围可以对应完全不同的更新轨迹。网络误差主要由每一步的增量分布决定，而不只由 $G_{\max}/G_{\min}$ 决定。}
\end{figure}

\begin{warningbox}
用一个“非线性系数”压缩整条曲线时，要明确拟合模型、脉冲数和归一化方式。不同论文采用的参数定义并不等价，不能把数字直接横向比较。
\end{warningbox}

\section{随机性必须分解：循环内、循环间与器件间}

至少区分三类量：
\begin{itemize}
\item \textbf{写入步长随机性：}同一状态附近每个脉冲产生的 $\Delta G$ 波动；
\item \textbf{循环间波动（C2C）：}同一器件重复完整增强--抑制周期后的差异；
\item \textbf{器件间差异（D2D）：}阵列中不同像素的动态范围、阈值、线性度和保持差异。
\end{itemize}
网络仿真应从实测分布采样，而不是只给平均值加一个任意高斯噪声。若误差随状态、脉冲方向或时间变化，使用常方差噪声会给出过度乐观结果。

\begin{figure}[H]
\centering
\begin{tikzpicture}[x=0.8cm,y=0.65cm]
\draw[->] (0,0)--(10.5,0) node[right]{归一化电导};
\draw[->] (0,0)--(0,6.3) node[above]{概率密度};
\foreach \x/\h in {1/0.7,2/1.3,3/2.5,4/4.8,5/5.5,6/4.2,7/2.2,8/1.0,9/0.4}{\fill[titleblue!65] (\x-0.35,0) rectangle (\x+0.35,\h);}
\draw[very thick,deepred] plot[smooth] coordinates {(1,0.2)(2,0.8)(3,2.0)(4,4.0)(5,5.7)(6,4.6)(7,2.5)(8,1.0)(9,0.3)};
\node[align=left,anchor=west] at (6.6,5.2){直方图：全部器件\\曲线：拟合仅作摘要};
\end{tikzpicture}
\caption{阵列统计必须显示原始分布、样本数和失效规则。只报告均值与标准差会隐藏偏态、双峰分布和坏点群。}
\end{figure}

\section{权重编码：差分对并不会免费消除非理想性}

正负权重常用差分电导表示：
\[
 w_{ij}=\alpha\left(G^+_{ij}-G^-_{ij}\right).
\]
它能表示符号并部分抵消共模漂移，却将面积、写入次数和静态电流近似加倍。若 $G^+$ 与 $G^-$ 的更新不相关，权重噪声为
\[
\sigma_w^2=\alpha^2\left(\sigma_{G^+}^2+\sigma_{G^-}^2\right).
\]
因此差分编码不是“自动降噪”，其收益取决于误差的相关性。

\section{交叉阵列乘加：理想公式与真实电路之间}

理想列电流满足
\[
I_j=\sum_i G_{ij}V_i.
\]
真实阵列还包含线阻、接触电阻、半选泄漏、器件非线性、读扰动和 ADC 量化。尺寸增大后，远端单元看到的实际电压小于输入端电压，导致位置相关误差。网络映射必须说明使用的是理想矩阵、等效电路仿真，还是实测阵列输出。

\begin{figure}[H]
\centering
\begin{tikzpicture}[node distance=8mm,box/.style={draw=deepblue,rounded corners=2mm,fill=softblue,minimum width=27mm,minimum height=11mm,align=center},err/.style={draw=deepred,rounded corners=2mm,fill=softred,minimum width=27mm,minimum height=11mm,align=center},arr/.style={-{Latex},thick}]
\node[box] (data) {输入编码\\DAC/脉冲};
\node[box,right=of data] (array) {器件阵列\\$\mathbf I=\mathbf G\mathbf V$};
\node[box,right=of array] (read) {读出\\TIA/ADC};
\node[box,right=of read] (task) {软件或硬件\\后处理};
\draw[arr] (data)--(array); \draw[arr] (array)--(read); \draw[arr] (read)--(task);
\node[err,below=of data] (e1) {量化与驱动能耗};
\node[err,below=of array] (e2) {线阻、D2D、漂移};
\node[err,below=of read] (e3) {噪声、带宽、饱和};
\draw[dashed] (e1)--(data); \draw[dashed] (e2)--(array); \draw[dashed] (e3)--(read);
\end{tikzpicture}
\caption{从权重到任务输出的数据路径。论文若只计算阵列内的 $GV$，却把外围电路排除在延迟和能耗之外，应明确限定结论边界。}
\end{figure}

\section{训练方式的四个层级}

\begin{longtable}{>{\bfseries\color{titleblue}}p{0.19\textwidth}p{0.34\textwidth}p{0.37\textwidth}}
\toprule
层级 & 实际做了什么 & 必须公开的内容\\
\midrule
理想软件基线 & 连续高精度权重训练 & 数据集、网络、训练轮数、随机种子\\
器件感知训练 & 在前向或更新中加入拟合非理想性 & 参数分布来源、相关性与采样方法\\
硬件在环 & 中间量由真实器件或阵列读出 & 硬件参与步骤、采样次数、校准和接口\\
原位训练 & 权重更新直接发生在阵列 & 写验证策略、误差反馈、耗时与耐久损耗\\
\bottomrule
\end{longtable}

\section{准确率下降要能归因，而不是只给一个数字}

建议进行逐项消融：先加入更新非线性，再加入 C2C、D2D、保持漂移、读噪声和量化。若同时加入全部误差，只能得到“下降了多少”，不能知道下一步应优化材料、脉冲、阵列还是算法。

\begin{examplebox}
\textbf{例 16.1（把识别率差距拆开）}\quad 理想软件准确率为 $98.1\%$；加入实测更新非线性后为 $96.7\%$；再加入 D2D 后为 $94.2\%$；加入 24 h 漂移后为 $92.8\%$。此时最优先的器件改进对象不是动态范围，而是器件间参数分散和保持漂移。
\end{examplebox}

\section{能耗与延迟：必须画清系统边界}

单次写脉冲能量可近似为
\[
E_{\mathrm{pulse}}=\int V(t)I(t)\,\dd t,
\]
但任务能耗还包括脉冲发生、地址译码、写验证、读出、ADC/DAC、数据搬运和校准。延迟同样不只是器件响应时间，而是
\[
t_{\mathrm{task}}=t_{\mathrm{program}}+t_{\mathrm{settle}}+t_{\mathrm{read}}+t_{\mathrm{convert}}+t_{\mathrm{compute}}.
\]

\begin{checkbox}
\begin{itemize}
\tickitem 已把增强、抑制、保持和读噪声转换为可复现模型；
\tickitem 网络仿真从实测分布采样，而不是只使用平均曲线；
\tickitem 已区分理想软件、器件感知、硬件在环与原位训练；
\tickitem 准确率下降有逐项误差归因；
\tickitem 能耗与延迟包含外围电路和校准边界；
\tickitem 阵列规模、有效良率和任务容差同时报告。
\end{itemize}
\end{checkbox}

\begin{conclusionbox}
器件曲线只有在被转换为权重模型、误差分布、电路约束和任务协议后，才具有系统意义。高质量工作不应只回答“软件准确率是多少”，而应回答“哪些实测非理想造成了多少损失，哪一种材料或电路改进最值得做”。
\end{conclusionbox}

\chapter{光谱课题怎样选问题}

光谱实验最容易陷入“测到峰、认出峰、画出峰移”的低水平循环。可投稿的光谱课题必须把谱线变化连接到一个可证伪的物理问题，并通过功率、温度、栅压、时间、偏振或空间维度排除竞争解释。

\begin{mainline}
一个可执行的光谱课题应写成
\[
\boxed{\text{对象}}+\boxed{\text{可控变量}}+\boxed{\text{观测量}}+\boxed{\text{竞争机制}}.
\]
例如：“单层 WSe$_2$ 低能 PL 峰在栅压下是否遵循带电激子人口变化，而不是固定缺陷态的热激活？”这比“研究 WSe$_2$ 的 PL 性质”更接近真实课题。
\end{mainline}

\section{先决定要回答的是能量、人口、动力学还是空间分布}

\begin{longtable}{>{\bfseries\color{titleblue}}p{0.18\textwidth}p{0.30\textwidth}p{0.42\textwidth}}
\toprule
问题类型 & 主要观测量 & 常见实验组合\\
\midrule
能级与耦合 & 峰位、分裂、反交叉 & 反射/吸收、PL、栅压、磁场、温度\\
人口与复合 & 峰面积、量子产率、功率律 & PL、PLE、功率、栅压、温度\\
动力学 & 寿命、形成与转移时间 & TRPL、泵浦--探测、全局动力学拟合\\
声子与晶格 & Raman 位移、线宽、强度 & 偏振、温度、应变、低频 Raman\\
对称性 & SHG 极图、相位、畴分布 & 偏振 SHG、相位干涉、外场和空间成像\\
空间非均匀 & 峰位/强度/寿命地图 & 共焦 mapping、nano-PL、AFM 联测\\
\bottomrule
\end{longtable}

\section{最小实验矩阵来自竞争机制}

\begin{table}[H]
\centering\small
\begin{tabularx}{\textwidth}{>{\bfseries\color{titleblue}}p{0.25\textwidth}X X}
\toprule
待区分问题 & 首选变量 & 决定性补充证据\\
\midrule
三激子还是缺陷峰 & 栅压、功率 & 温度、寿命、空间定位\\
PL 淬灭是否为电荷转移 & TRPL/瞬态吸收 & 界面依赖、对照层、温度\\
Raman 峰移来自应变还是掺杂 & 多模相关移动 & 独立门控或应变标定\\
线宽增大来自温度还是无序 & 功率/温度交叉扫描 & 激光热校准、空间统计\\
SHG 变化来自对称性还是干涉 & 偏振与相位 & 厚度、波长、参考晶体\\
\bottomrule
\end{tabularx}
\end{table}

\begin{translationbox}
“多测几个功率”不是实验矩阵。矩阵意味着每个变量都对应一个竞争机制的可检验预测，并提前规定什么结果支持、反驳或仍无法区分该机制。
\end{translationbox}

\section{数据处理链必须可回溯}

\begin{figure}[H]
\centering
\begin{tikzpicture}[node distance=6mm,box/.style={draw=deepblue,rounded corners=2mm,fill=softblue,minimum width=25mm,minimum height=10mm,align=center},arr/.style={-{Latex},thick}]
\node[box] (raw) {原始计数\\暗背景/元数据};
\node[box,right=of raw] (cal) {仪器校正\\响应/波数/IRF};
\node[box,right=of cal] (fit) {模型拟合\\峰位/面积/线宽};
\node[box,right=of fit] (phys) {物理参数\\速率/耦合/应变};
\node[box,below=of phys] (claim) {机制主张\\与竞争解释比较};
\draw[arr] (raw)--(cal); \draw[arr] (cal)--(fit); \draw[arr] (fit)--(phys); \draw[arr] (phys)--(claim);
\draw[arr,dashed] (claim.west)-| node[below,sloped]{可回查} (raw.south);
\end{tikzpicture}
\caption{光谱证据链。最终机制结论必须能够回查到原始计数、校准文件、拟合残差和参数不确定度。}
\end{figure}

\section{稳态、时间分辨与空间成像怎样分工}

稳态光谱回答“哪些辐射通道在积分时间内贡献了多少”；时间分辨实验回答“人口怎样产生、转移和衰减”；空间成像回答“现象是否与界面、褶皱、边缘、畴或污染共定位”。三者互补，不能用一张高信噪比稳态 PL 代替动力学证据。

\section{校准不是附录，而是课题的一部分}

最低记录包括：样品面功率、光斑尺寸、积分时间、物镜、光栅、狭缝、偏振元件、环境、温度、光谱响应校正、波长/波数标定和重复扫描漂移。对功率或温度依赖实验，还应测试激光加热与光致掺杂。

\begin{warningbox}
同一条谱线随功率红移，可能来自激子相互作用，也可能只是样品升温；PL 增强可能来自辐射速率提高，也可能来自吸收增加或收集效率变化。未经校准的强度和峰移不能直接上升为材料机制。
\end{warningbox}

\section{怎样判断课题是否具有新增量}

把已有文献结论拆成四层：已知现象、已知机制、已有控制手段、尚未解释的矛盾。新增量通常来自：
\begin{itemize}
\item 在更干净或更复杂体系中推翻旧峰归属；
\item 用新的可控变量区分此前混淆的机制；
\item 把稳态现象推进到时间或空间分辨；
\item 建立可预测的定量模型，而非只报告峰移趋势；
\item 将光谱指纹与器件状态或相变一一对应。
\end{itemize}

\section{光谱论文的最小主图包}

\begin{longtable}{>{\bfseries\color{titleblue}}p{0.13\textwidth}p{0.46\textwidth}p{0.31\textwidth}}
\toprule
图组 & 推荐内容 & 回答的问题\\
\midrule
Fig. 1 & 样品、结构、基础 Raman/PL 与校准 & 样品是谁、测量是否可信\\
Fig. 2 & 核心谱学现象及可重复统计 & 新现象是否真实\\
Fig. 3 & 功率/温度/栅压/偏振等变量依赖 & 哪些机制预测被满足\\
Fig. 4 & 时间或空间分辨、关键对照 & 能否排除竞争机制\\
Fig. 5 & 定量模型、器件关联或外场操控 & 物理图景是否可预测\\
SI & 原始谱、校准、拟合残差、更多样品 & 结果能否复查\\
\bottomrule
\end{longtable}

\begin{checkbox}
\begin{itemize}
\tickitem 课题已写成“对象+变量+观测量+竞争机制”；
\tickitem 每个变量都有明确可证伪预测；
\tickitem 强度、峰位和线宽经过相应仪器校正；
\tickitem 核心结论至少有两个互补维度支持；
\tickitem 原始谱、拟合残差和参数不确定度可追溯；
\tickitem 已说明仍不能由当前实验区分的解释。
\end{itemize}
\end{checkbox}

\begin{conclusionbox}
高质量光谱课题不是“峰越多越好”，而是用最少但互补的实验维度把峰的物理来源限定到一个可检验模型。选题时先列竞争机制，再决定仪器和变量，通常比先测大量光谱再寻找故事更有效。
\end{conclusionbox}

\chapter{PL 峰归属与激子物理}

二维半导体 PL 中，中性激子、带电激子、双激子、暗态、缺陷束缚态和应变局域态可能在几十毫电子伏范围内重叠。峰拟合只能回答“数学上需要几个分量”，不能单独回答“每个分量是什么”。

\section{PL 峰不是单粒子带隙}

吸收形成的电子--空穴对受库仑作用束缚，光学跃迁能量近似为
\[
E_{X}=E_{\mathrm g}^{\mathrm{QP}}-E_{\mathrm b}^{X}.
\]
因此 PL 峰给出的是激子复合能，而非准粒子带隙。衬底、封装和载流子密度会同时改变准粒子带隙和束缚能，两者可能部分抵消，导致峰位变化远小于各自变化。

\section{常见发光态及其最强判据}

\begin{longtable}{>{\bfseries\color{titleblue}}p{0.17\textwidth}p{0.29\textwidth}p{0.44\textwidth}}
\toprule
候选态 & 典型现象 & 更有区分力的证据\\
\midrule
中性激子 $X^0$ & 主发光峰、近带边 & 栅压去掺杂增强、偏振与反射共振对应\\
带电激子 $X^{\pm}$ & 比 $X^0$ 低能、随载流子变化 & 可逆栅压人口转移、质量作用关系\\
双激子 $XX$ & 高激发下出现低能峰 & 超线性功率依赖、关联/时间动力学\\
缺陷束缚态 & 宽峰、低能、样品差异大 & 空间定位、热激活、缺陷处理前后对照\\
应变局域态 & 褶皱/纳米结构处低能发光 & AFM/应变图与 nano-PL 共定位\\
暗激子 & 低温或磁场下激活 & 磁场、温度、面内场或声子辅助选择定则\\
层间激子 & 异质结低能、寿命较长 & 层间偶极 Stark 移、界面依赖、时间分辨\\
\bottomrule
\end{longtable}

\section{多峰拟合：从线型到参数约束}

常用模型为
\[
I(E)=B(E)+\sum_{i=1}^{m}A_i L_i(E;E_i,\Gamma_i),
\]
其中 $A_i$ 应定义为积分面积而非峰高。均匀寿命展宽常近似 Lorentzian，静态无序常导致 Gaussian 分量；两者卷积为 Voigt。真实谱还需卷积仪器响应。

\begin{figure}[H]
\centering
\begin{tikzpicture}
\begin{axis}[width=0.86\textwidth,height=6.3cm,xlabel={能量 / eV},ylabel={PL 强度 / a.u.},xmin=1.68,xmax=1.98,ymin=0,ymax=1.45,grid=both,samples=220,legend style={at={(0.02,0.98)},anchor=north west,draw=none,fill=white}]
\addplot[very thick,black,domain=1.68:1.98]{0.05+0.92*exp(-((x-1.91)/0.025)^2)+0.55*exp(-((x-1.865)/0.030)^2)+0.24*exp(-((x-1.78)/0.055)^2)};
\addlegendentry{总谱（教学模拟）}
\addplot[thick,titleblue,dashed,domain=1.68:1.98]{0.92*exp(-((x-1.91)/0.025)^2)};
\addlegendentry{$X^0$}
\addplot[thick,deepgreen,dashed,domain=1.68:1.98]{0.55*exp(-((x-1.865)/0.030)^2)};
\addlegendentry{$X^-$}
\addplot[thick,gold,dashed,domain=1.68:1.98]{0.24*exp(-((x-1.78)/0.055)^2)};
\addlegendentry{局域/缺陷候选峰}
\end{axis}
\end{tikzpicture}
\caption{教学模拟：谱线数学分解。即使总拟合非常好，低能宽峰仍可能对应缺陷、应变局域态或多个未分辨状态，峰名必须由外部变量决定。}
\end{figure}

\section{避免过拟合：峰数不是越多越好}

至少同时检查：残差是否具有系统结构、参数置信区间、峰位和线宽是否跨数据集连续、增加峰后 AIC/BIC 是否显著改善，以及峰是否在多个样品或控制变量中复现。对温度或栅压序列，优先全局拟合：让共享参数跨谱线共同约束，而不是每条谱独立自由拟合。

\begin{warningbox}
若新增峰的线宽接近仪器分辨率、面积与背景高度强相关，或峰位在相邻条件下随机跳动，该峰通常不具备稳定物理身份。一个低残差图不能替代参数可辨识性分析。
\end{warningbox}

\section{栅压：区分中性激子与带电激子的首选变量}

改变载流子密度时，$X^0$ 与 $X^-$ 常发生可逆人口转移。简化质量作用关系可写成
\[
\frac{n_X n_e}{n_{X^-}}=K(T),
\]
其中 $K(T)$ 取决于有效质量、温度和三激子束缚能。实验上应同时报告栅漏电、阈值漂移和迟滞，避免把不可逆陷阱充电误判为平衡人口转移。

\section{功率、温度和寿命：低能峰的联合判据}

\begin{itemize}
\item 双激子候选峰通常在低密度区呈超线性增长，但高密度下会受湮灭和饱和影响；
\item 缺陷/局域态常表现出饱和、热淬灭和更强空间非均匀；
\item 暗态可能随温度或磁场激活，不能只由低能位置命名；
\item 寿命更长不自动意味着“缺陷”，也可能来自动量间接或层间态。
\end{itemize}

\begin{auditbox}
\textbf{案例一：MoS$_2$ 三激子与衬底效应。} Drüppel 等的开放论文指出，负三激子可包含多个谷内/谷间激发，且衬底屏蔽会同时重整化带隙、激子与三激子束缚能。教材上的“一个 $X^-$ 峰”只是最小模型；高分辨实验中发现分裂时，应先考虑多体态与介电环境，而不是机械增加缺陷峰。\\
论文：Nature Communications 8, 2117 (2017)，CC BY 4.0，DOI: 10.1038/s41467-017-02286-6。
\end{auditbox}

\begin{auditbox}
\textbf{案例二：纳米褶皱诱导的局域激子。} Yanev 等将 AFM 形貌和应变模型与 nano-PL 超光谱图逐点关联，证明低能发光与细褶皱处应变局域高度共定位，并在低温观察到反聚束发光。这个证据链比“低能峰+功率饱和”强得多，因为它把结构、空间和光子统计连接起来。\\
论文：Nature Communications 15, 1543 (2024)，CC BY 4.0，DOI: 10.1038/s41467-024-45936-2。
\end{auditbox}

\section{PL 峰归属的证据等级}

\begin{table}[H]
\centering\small
\begin{tabularx}{\textwidth}{>{\bfseries\color{titleblue}}p{0.16\textwidth}X X}
\toprule
等级 & 可以合理表述 & 仍缺少什么\\
\midrule
峰位相容 & “低能峰与三激子能量范围相容” & 栅压、温度或人口证据\\
功率/温度支持 & “表现出局域态常见的饱和与热淬灭” & 空间或缺陷控制\\
可逆外场操控 & “栅压驱动 $X^0/X^-$ 人口转移” & 平衡性、迟滞与重复性\\
空间/时间共定位 & “峰与褶皱/界面及特定寿命共定位” & 替代结构与多样品验证\\
直接选择定则/统计 & 磁场、偏振或反聚束锁定状态性质 & 仪器响应与完整模型\\
\bottomrule
\end{tabularx}
\end{table}

\begin{checkbox}
\begin{itemize}
\tickitem 每个峰的面积、峰位、线宽和不确定度均已报告；
\tickitem 峰数经过残差与模型比较，而不是凭经验固定；
\tickitem 峰归属至少由栅压、温度、功率、时间或空间中的两维支持；
\tickitem 已区分可逆人口转移与不可逆陷阱充电；
\tickitem 低能峰未被默认命名为“缺陷峰”；
\tickitem 原始谱、背景模型和全局拟合约束可复查。
\end{itemize}
\end{checkbox}

\begin{conclusionbox}
PL 峰归属是一个模型选择问题。拟合负责提取可比较参数，栅压、功率、温度、时间、偏振和空间证据负责赋予物理身份。真正可靠的归属应在多个变量下给出同一套可预测图景。
\end{conclusionbox}

\chapter{温度、功率和时间分辨 PL}

稳态 PL 强度是产生、弛豫、辐射复合、非辐射复合和收集效率共同作用的结果。功率、温度和时间分辨实验的价值，在于把这些过程拆开；它们不能被当成三组彼此独立的“补充曲线”。

\section{功率依赖：从速率方程理解次线性}

最小人口模型为
\[
\frac{\dd n}{\dd t}=G-\frac{n}{\tau}-k_{\mathrm A}n^2,
\]
其中 $G\propto P$，$k_{\mathrm A}$ 表示激子--激子湮灭或其他双体损失。低密度时 $n\approx G\tau$，PL 近线性；高密度时双体项增强，PL 变为次线性。经验功率律
\[
I_{\mathrm{PL}}\propto P^{\alpha}
\]
只是局部斜率摘要，$\alpha<1$ 并不唯一证明缺陷或湮灭。

\begin{figure}[H]
\centering
\begin{tikzpicture}
\begin{loglogaxis}[width=0.82\textwidth,height=6cm,xlabel={激发功率 / a.u.},ylabel={积分 PL / a.u.},xmin=0.01,xmax=100,ymin=0.005,ymax=30,grid=both,legend style={at={(0.03,0.97)},anchor=north west,draw=none}]
\addplot[very thick,titleblue,domain=0.01:100,samples=100]{x/(1+0.35*x)^0.45};
\addlegendentry{含双体损失}
\addplot[thick,deepgreen,dashed,domain=0.01:100]{x};
\addlegendentry{理想线性}
\addplot[thick,gold,dashdotted,domain=0.01:100]{0.8*x^0.55};
\addlegendentry{局域态饱和型}
\end{loglogaxis}
\end{tikzpicture}
\caption{教学模拟：不同机制可在有限功率窗口内给出相似斜率。应比较完整速率模型、谱形变化和重复升降功率扫描。}
\end{figure}

\section{温度峰移：Varshni 与声子模型只是经验参数化}

常用 Varshni 形式为
\[
E(T)=E(0)-\frac{\alpha T^2}{T+\beta},
\]
也可使用 Bose--Einstein 型声子耦合模型。拟合参数在有限温区内高度相关，不应无误差地解释为唯一声子能量。对异质结或局域态，还要考虑热人口转移和不同峰重叠造成的表观峰移。

\section{线宽随温度：把均匀展宽与静态无序分开}

一种常见经验式是
\[
\Gamma(T)=\Gamma_0+\gamma_{\mathrm{ac}}T+rac{\Gamma_{\mathrm{LO}}}{\exp(E_{\mathrm{LO}}/k_{\mathrm B}T)-1}.
\]
$\Gamma_0$ 包含静态无序和仪器贡献，线性项近似声学声子，最后一项表示光学声子散射。若不同温度下峰之间发生人口转移，单峰线宽模型会失效，应使用全局多峰拟合。

\section{热淬灭：Arrhenius 激活能不自动等于束缚能}

常用模型
\[
I(T)=\frac{I_0}{1+C\exp[-E_a/(k_{\mathrm B}T)]}
\]
假设单一热激活非辐射通道。实际数据常需要多个通道或温度依赖吸收。$E_a$ 只能称为“有效激活能”，除非有独立证据把它对应到某个束缚能或势垒。

\begin{warningbox}
温度升高导致 PL 下降并不能单独说明激子解离。非辐射缺陷、暗亮态人口转换、层间转移和收集效率变化都可能产生热淬灭。必须观察谱峰人口、寿命或吸收的同步变化。
\end{warningbox}

\section{TRPL：测到的是样品响应与仪器响应的卷积}

时间相关单光子计数测得
\[
I_{\mathrm{meas}}(t)=\bigl[\mathrm{IRF}*I_{\mathrm{true}}\bigr](t)+B.
\]
当寿命接近 IRF 宽度时，直接在半对数图上拟合指数会严重偏差。应使用重卷积拟合，并把时间零点、背景和 IRF 漂移纳入模型。

\begin{figure}[H]
\centering
\begin{tikzpicture}
\begin{axis}[width=0.84\textwidth,height=6.2cm,xlabel={时间 / ns},ylabel={归一化计数},xmin=-0.15,xmax=3.5,ymin=0,ymax=1.1,grid=both,legend style={at={(0.98,0.97)},anchor=north east,draw=none},samples=200]
\addplot[very thick,gray,domain=-0.15:1]{exp(-((x-0.04)/0.12)^2)};
\addlegendentry{IRF}
\addplot[very thick,titleblue,domain=0:3.5]{0.72*exp(-x/0.28)+0.28*exp(-x/1.35)};
\addlegendentry{真实双指数模型}
\addplot[very thick,deepred,dashed,domain=-0.1:3.5]{0.62*exp(-max(x,0)/0.36)+0.28*exp(-max(x,0)/1.45)+0.10*exp(-((x-0.04)/0.15)^2)};
\addlegendentry{卷积后测量外形}
\end{axis}
\end{tikzpicture}
\caption{教学模拟：快寿命接近 IRF 时，测量曲线前沿由仪器主导。报告寿命时必须给出 IRF 和重卷积残差。}
\end{figure}

\section{多指数寿命：数学分量不等于独立微观通道}

经验式
\[
I(t)=\sum_i A_i\exp(-t/\tau_i)
\]
能摘要多时间尺度，但 $A_i$ 和 $\tau_i$ 通常强相关。平均寿命的定义也有多种：振幅加权与积分加权不能混用。机制分析应结合谱分辨寿命、温度、功率或泵浦--探测，而不是把每个指数分量直接命名为一种缺陷。

\section{寿命与量子产率共同给出速率}

在单一有效通道近似下
\[
\eta=\frac{k_r}{k_r+k_{nr}},\qquad \tau=\frac{1}{k_r+k_{nr}},
\]
于是
\[
k_r=\frac{\eta}{\tau},\qquad k_{nr}=\frac{1-\eta}{\tau}.
\]
仅有 PL 增强不能判断是 $k_r$ 增大还是 $k_{nr}$ 降低；仅有寿命缩短也可能来自辐射速率提高或非辐射损失增强。

\begin{auditbox}
\textbf{开放论文：化学处理后的 PL、Raman 与 TRPL 联合证据。} Li 等比较 H-TFSI 和 Li-TFSI 处理的 MoS$_2$/WS$_2$，不仅展示 PL 增强和蓝移，还用 Raman、XPS、TRPL、泵浦--探测和扩散测量排除“简单 p 掺杂”解释。Li-TFSI 处理样品 PL 更强但寿命更短，结合缺陷瞬态特征消失，说明辐射复合效率提高而非陷阱导致的长寿命增强。\\
论文：Nature Communications 12, 6044 (2021)，CC BY 4.0，DOI: 10.1038/s41467-021-26340-6。
\end{auditbox}

\section{一套可执行的联合实验协议}

\begin{enumerate}
\item 在低功率下标定吸收线性与无加热区间；
\item 进行升功率与降功率扫描，检查光致掺杂和不可逆漂移；
\item 选择若干功率做完整谱分解，而非只跟踪单点峰高；
\item 在相同位置和功率密度下进行温度序列；
\item 对关键峰做谱分辨 TRPL，并每次记录 IRF；
\item 用全局速率模型同时拟合功率、温度和时间数据；
\item 在多个样品和位置上报告参数分布。
\end{enumerate}

\begin{checkbox}
\begin{itemize}
\tickitem 功率律拟合限定了功率窗口并报告升降扫描；
\tickitem 已排查激光加热、光致掺杂和探测器饱和；
\tickitem 温度拟合参数含置信区间，激活能未被过度命名；
\tickitem TRPL 使用实测 IRF 重卷积并显示残差；
\tickitem 多指数分量未被直接等同于多个微观缺陷；
\tickitem PLQY 与寿命联合用于区分辐射和非辐射速率。
\end{itemize}
\end{checkbox}

\begin{conclusionbox}
功率、温度和时间分辨 PL 应由同一个人口动力学模型连接。高质量工作不仅报告“峰如何变”，还要说明产生、复合、转移和损失速率怎样随外部变量变化，并明确哪些参数仍只是有效量。
\end{conclusionbox}

\chapter{Raman、应变、掺杂与层间耦合}

Raman 峰位是声子自能、应变、载流子、温度、层间耦合和共振条件共同作用的结果。把所有红移称为“拉伸应变”、把所有线宽变化称为“缺陷增加”，是二维材料课题中最常见的过度解释之一。

\section{一张 Raman 谱至少包含四类信息}

\begin{itemize}
\item \textbf{峰位：}力常数、应变、载流子与温度；
\item \textbf{线宽：}声子寿命、电子--声子耦合、无序及仪器分辨率；
\item \textbf{强度：}Raman 张量、偏振、共振、干涉与聚焦条件；
\item \textbf{线型：}离散声子与连续电子激发耦合时可能出现 Fano 非对称。
\end{itemize}
因此峰位以外的维度不能被丢弃，尤其在掺杂和共振 Raman 中。

\section{应变与掺杂：使用多模相关移动}

设两个 Raman 模的位移为
\[
\begin{pmatrix}\Delta\omega_1\\\Delta\omega_2\end{pmatrix}
=
\begin{pmatrix}a_1&b_1\\a_2&b_2\end{pmatrix}
\begin{pmatrix}\varepsilon\\\Delta n\end{pmatrix},
\]
其中 $a_i$ 为应变系数，$b_i$ 为载流子系数。只有当两个响应向量不共线时，才能从相关图中分离 $\varepsilon$ 与 $\Delta n$。系数必须在相同材料、层数、温度和衬底条件下标定。

\begin{figure}[H]
\centering
\begin{tikzpicture}
\begin{axis}[width=0.76\textwidth,height=6.5cm,xlabel={$\Delta\omega_{\mathrm{in}}$ / cm$^{-1}$},ylabel={$\Delta\omega_{\mathrm{out}}$ / cm$^{-1}$},xmin=-4,xmax=4,ymin=-5,ymax=5,axis lines=middle,grid=both]
\addplot[very thick,titleblue,domain=-3.5:3.5]{0.35*x};
\node[titleblue] at (axis cs:2.5,1.5){掺杂方向};
\addplot[very thick,deepred,domain=-3.2:3.2]{1.25*x};
\node[deepred] at (axis cs:-2.4,-3.7){应变方向};
\addplot[only marks,mark=*,deepgreen] coordinates {(-2.2,-2.1)(-1.5,-1.2)(-0.8,-0.5)(0.2,0.5)(1.1,1.4)(2.1,2.2)};
\end{axis}
\end{tikzpicture}
\caption{教学示意：二维相关图把多模峰移投影到应变与掺杂基向量。若基向量近共线或系数误差很大，分解会病态。}
\end{figure}

\section{温度与激光加热必须单独标定}

激光功率改变 Raman 峰位时，应先用低功率线性区或 Stokes/anti-Stokes 强度比估算局部温升。即便样品台温度恒定，吸收、界面热阻和光斑大小仍会使局部温度不同。mapping 中边缘和气泡常具有不同散热条件，可能伪装成应变分布。

\begin{warningbox}
“峰位随激光功率红移”首先证明的是光照改变了声子能量，不自动证明光致应变。若没有独立温度标定、重复扫描和功率恢复测试，最保守解释通常是局部加热。
\end{warningbox}

\section{偏振 Raman：晶轴判定必须包含光路响应}

Raman 强度满足
\[
I\propto \left|\mathbf e_s^{\mathsf T}\mathbf R\mathbf e_i\right|^2,
\]
其中 $\mathbf R$ 为 Raman 张量。实际极图还受到物镜、反射镜、光栅和 CCD 偏振响应影响。应使用参考样品或系统传递矩阵校正，并同时报告入射与分析偏振几何。

\section{低频 Raman：剪切模与层间呼吸模}

多层体系可用线性链模型近似。对 $N$ 层、相邻层耦合常数 $\kappa$、面密度 $\mu$，模频率具有
\[
\omega_j\propto \sqrt{\frac{2\kappa}{\mu}\left[1-\cos\left(\frac{j\pi}{N}\right)\right]},\qquad j=1,\ldots,N-1.
\]
剪切模主要反映面内相对滑移，呼吸模反映面外相对运动；堆垛、扭转、压力和界面污染都会改变频率、选择定则与线宽。

\begin{figure}[H]
\centering
\begin{tikzpicture}
\begin{axis}[width=0.82\textwidth,height=6.2cm,xlabel={层数 $N$},ylabel={归一化低频模频率},xmin=2,xmax=8,ymin=0,ymax=1.1,xtick={2,3,4,5,6,7,8},grid=both,legend style={at={(0.98,0.03)},anchor=south east,draw=none}]
\addplot[very thick,titleblue,mark=*] coordinates {(2,0.71)(3,0.87)(4,0.92)(5,0.95)(6,0.97)(7,0.98)(8,0.99)};
\addlegendentry{最高剪切分支}
\addplot[very thick,deepgreen,mark=square*] coordinates {(2,0.71)(3,0.50)(4,0.38)(5,0.31)(6,0.26)(7,0.22)(8,0.20)};
\addlegendentry{最低呼吸/剪切分支}
\end{axis}
\end{tikzpicture}
\caption{线性链教学模拟：层数增加会产生多个分支。真实峰的可见性还由堆垛点群和偏振选择定则决定。}
\end{figure}

\section{共振 Raman：强度峰不一定表示声子数量增加}

当激光能量接近激子共振时，Raman 截面可增强几个数量级，并出现多声子或原本弱的模式。此时不同激发能下强度不能简单归一化比较，必须同时测量反射/PL/PLE 或建立激发能扫描。峰强变化可能来自入射共振、出射共振或激子--声子耦合改变。

\begin{auditbox}
\textbf{开放论文：四种单层 TMD 的 Raman scattering excitation。} Zinkiewicz 等在 5 K 下对 hBN 封装的 MoS$_2$、MoSe$_2$、WS$_2$ 和 WSe$_2$ 扫描激发能，在固定探测能下构建 Raman scattering excitation 光谱，显示声子特征与激子共振之间的材料依赖关系。这个案例说明：Raman 强度必须放回激发能和激子结构中解释，而不能只在单一激光波长下比较。\\
论文：npj 2D Materials and Applications 8, 2 (2024)，CC BY 4.0，DOI: 10.1038/s41699-023-00438-5。
\end{auditbox}

\section{缺陷、无序与 Fano 线型}

缺陷可能放松动量守恒、激活边界声子并增加非均匀展宽，但峰宽增加也可能来自载流子、温度或相邻峰重叠。Fano 线型可写为
\[
I(\omega)=I_0\frac{(q+\epsilon)^2}{1+\epsilon^2},\qquad \epsilon=\frac{\omega-\omega_0}{\Gamma/2},
\]
其中 $q$ 描述离散声子与连续背景的干涉。用对称 Lorentzian 强行拟合会造成假峰移和假线宽。

\section{Raman mapping：每个像素都应保留拟合质量}

mapping 不应只输出一张平滑颜色图。每个像素至少保存峰位、线宽、面积、背景、信噪比和拟合残差；同时记录步长、光斑、积分时间和插值方法。对边缘、气泡和异质结，应叠加 AFM 或显微图进行空间共定位。

\section{一套应变--掺杂--层间耦合实验矩阵}

\begin{table}[H]
\centering\small
\begin{tabularx}{\textwidth}{>{\bfseries\color{titleblue}}p{0.20\textwidth}X X}
\toprule
变量 & 主要作用 & 必要对照\\
\midrule
单轴/双轴应变 & 模分裂、相关峰移 & 独立应变标尺、偏振几何\\
栅压 & 载流子与电子--声子耦合 & 漏电、迟滞、接触和温升\\
温度/激光功率 & 非谐性与局部加热 & 低功率恢复、anti-Stokes\\
层数/堆垛 & 低频模与选择定则 & AFM、SHG 或结构证据\\
扭转角 & 层间耦合、区折叠与莫尔声子 & 角度标定、多个区域\\
\bottomrule
\end{tabularx}
\end{table}

\begin{checkbox}
\begin{itemize}
\tickitem Raman 峰位、线宽、面积和线型均被检查；
\tickitem 应变与掺杂由多模相关移动或独立控制分离；
\tickitem 已标定激光加热和系统偏振响应；
\tickitem 低频模归属包含层数、堆垛和选择定则；
\tickitem mapping 同时保存像素级残差和质量掩膜；
\tickitem 共振 Raman 强度与激发能/激子结构联合解释。
\end{itemize}
\end{checkbox}

\begin{conclusionbox}
Raman 光谱不是单峰温度计或应变计，而是一组彼此耦合的声子观测量。最可信的课题通过多模、偏振、温度、栅压和低频信息建立过定约束，使应变、掺杂、加热和层间耦合不再只靠故事区分。
\end{conclusionbox}

\chapter*{第十六至第二十章参考文献}
\addcontentsline{toc}{chapter}{第十六至第二十章参考文献}
\begin{enumerate}[label={[\arabic*]},leftmargin=3.5em]
\setcounter{enumi}{45}
\item S. B. Hamid, J. A. C. Incorvia and S. Liu, ``Non-idealities in artificial synapses,'' \textit{Nature Reviews Physics} (2026).
\item D. Ielmini and H.-S. P. Wong, ``In-memory computing with resistive switching devices,'' \textit{Nature Electronics} \textbf{1}, 333--343 (2018).
\item M. Hu, J. P. Strachan, Z. Li, et al., ``Dot-product engine for neuromorphic computing: programming 1T1M crossbar to accelerate matrix-vector multiplication,'' DAC (2016).
\item M. Drüppel, T. Deilmann, P. Krüger, et al., ``Diversity of trion states and substrate effects in the optical properties of an MoS$_2$ monolayer,'' \textit{Nature Communications} \textbf{8}, 2117 (2017). DOI: 10.1038/s41467-017-02286-6. CC BY 4.0.
\item E. S. Yanev, T. P. Darlington, S. A. Ladyzhets, et al., ``Programmable nanowrinkle-induced room-temperature exciton localization in monolayer WSe$_2$,'' \textit{Nature Communications} \textbf{15}, 1543 (2024). DOI: 10.1038/s41467-024-45936-2. CC BY 4.0.
\item Z. Li, H. Bretscher, Y. Zhang, et al., ``Mechanistic insight into the chemical treatments of monolayer transition metal disulfides for photoluminescence enhancement,'' \textit{Nature Communications} \textbf{12}, 6044 (2021). DOI: 10.1038/s41467-021-26340-6. CC BY 4.0.
\item J. W. Christopher, B. B. Goldberg and A. K. Swan, ``Long tailed trions in monolayer MoS$_2$: Temperature dependent asymmetry and resulting red-shift of trion photoluminescence spectra,'' \textit{Scientific Reports} \textbf{7}, 14062 (2017).
\item M. Zinkiewicz, M. Grzeszczyk, T. Kazimierczuk, et al., ``Raman scattering excitation in monolayers of semiconducting transition metal dichalcogenides,'' \textit{npj 2D Materials and Applications} \textbf{8}, 2 (2024). DOI: 10.1038/s41699-023-00438-5. CC BY 4.0.
\item C. Lee, H. Yan, L. E. Brus, et al., ``Anomalous lattice vibrations of single- and few-layer MoS$_2$,'' \textit{ACS Nano} \textbf{4}, 2695--2700 (2010).
\item A. M. van der Zande, R. M. Huang, D. A. Chenet, et al., ``Grains and grain boundaries in highly crystalline monolayer molybdenum disulphide,'' \textit{Nature Materials} \textbf{12}, 554--561 (2013).
\end{enumerate}

\chapter*{本轮版本说明}
\addcontentsline{toc}{chapter}{本轮版本说明}
本续册完成第十六至第二十章：从器件非理想到网络映射、光谱课题设计、PL 峰归属、温度/功率/TRPL 以及 Raman 应变--掺杂--层间耦合。图形均为可复现的 TikZ/PGFPlots 教学图；三处开放论文审计采用 CC BY 4.0 论文的文字证据链，未缩小粘贴整张多面板原图，以避免在续册独立版中牺牲可读性。待与 v0.7 主册合并时，可再按“高 DPI 整页渲染--panel 精准裁剪--视觉复核--图源登记”流程加入真正不可替代的实验 panel。

下一轮将进入第二十一章“SHG、晶体对称性与铁电畴”、第二十二章“异质结、层间激子与莫尔光谱”、第二十三章“PL、Raman 与光电流联合判断异质结机制”、第二十四章“用光谱识别人工突触内部状态”和第二十五章“滑移铁电的 Raman、PL、SHG 与器件应用”。

\end{document}
